% Einheit 4 von 5 – Parameter aus Bedingungen bestimmen (bank/funktionsscharen-und-ortskurven/e4.jsonl, zusammenbau v0.1)
%% TODO Zeitmarke Einheit 4: Marken-Zeile nicht lesbar
%% TODO Zweigzeile Teil 1: was hier gelernt wird (Satz in Schülersprache aus dem Einheitstitel) – Einheit 4
\einheitenkopf[e4]{Einheit 4 von 5 · Parameter aus Bedingungen bestimmen}
\zweigzeile{Abitur LK}
\setcounter{aufgabe}{15}

%% TODO Titel: Ich-kann-Satz fehlt in der Bank (Platzhalter: Kettenname) – Nr. 16
%% TODO Anweisung: ein Satz über der ersten Teilaufgabe fehlt in der Bank – Nr. 16
\begin{aufgabe}{Parameter aus Bedingungen}
%% TODO \rechenplatz{n}: Schrittzahl des Grundfalls fehlt in der Bank (nur bei fester Schreibform, 2.3 b)
\begin{teile}
\teil Aufgabe: „Berechne den Hochpunkt von $f_3$ der Schar $f_a(x) = -x^2 + a x$.“ Vorwärts (Eigenschaft aus gegebenem Parameter) oder rückwärts (Parameter aus einer Bedingung)? \\ \kreuz{vorwärts} \kreuz{rückwärts}
\teil Berechne den Inhalt der Fläche, die der Graph von $f_a(x) = -x^2 + a x$ mit $a > 0$ und die x-Achse einschließen. $A =$ \_\_
\teil $H_k(x) = k \cdot (x - 2)^2 \cdot e^{x}$ mit $k > 0$ ist eine Stammfunktion von $h_k$. Der Graph von $h_k$ schließt über $[0; 2]$ mit der x-Achse ein Flächenstück ein. Bestimme $k$ so, dass es den Inhalt $1$ hat. \hfill (Abitur 2024 GK) \\ $k =$ \_\_
\teil Der Querschnitt eines Kübelfußes wird oben durch $p_c(x) = c \cdot (4 - x^2)$ mit $c > 0$ begrenzt, unten durch die x-Achse; 1 LE entspricht 5 cm. Der Fuß ist am Boden 20 cm breit und hat die Querschnittsfläche 400 cm². Bestimme $c$. \hfill (Abitur 2017 LK) \\ $c =$ \_\_
\teil Gegeben ist $f_k(x) = x^4 - 2k x^2$ mit $k > 0$. Beide Tiefpunkte haben die y-Koordinate $-9$. Bestimme $k$. \hfill (Abitur 2021 LK) \\ $k =$ \_\_
\teil Gegeben sind $f_k(x) = \sqrt{k x^2 + 100}$ mit $k > 0$ und $g(x) = \frac{1}{4}x^2 + 10$. Für welche $k$ haben die Graphen mehr als einen gemeinsamen Punkt? \hfill (Abitur 2018 LK)
\teil Eine Rampe folgt für $x \le 0$ dem Graphen von $f(x) = 1 + e^{2x}$ und danach der Parabel $p_a(x) = -a x^2 + b x + 2$ mit $a > 0$ ($x$ und $y$ in m). Der Übergang bei $x = 0$ soll knickfrei sein. Bestimme $b$. \hfill (Abitur 2018 LK) \\ $b =$ \_\_
\teil Das Profil einer Vase wird durch $h_k(x) = k - \frac{8}{x^2}$ mit $k > 3$ und $x > 0$ beschrieben; 1 LE entspricht 10 cm. Der Boden liegt auf der x-Achse, die Vase ist 30 cm hoch, und ihr Randradius ist doppelt so groß wie der Bodenradius. Bestimme $k$. \hfill (Abitur 2022 GK) \\ $k =$ \_\_
\teil Gegeben ist $f_a(x) = -a x^2 + 2x$ mit $a > 0$; ihr Hochpunkt ist $H(v \,|\, f_a(v))$. Die Punkte $O(0 | 0)$, $P(v | 0)$, $H$ und $Q(0 \,|\, 2v)$ bilden ein Viereck mit dem Inhalt $24$. Bestimme $a$; beginne mit einer Skizze. \hfill (Abitur 2022 LK) \\

\begin{ksys}[xmin=-1,xmax=6,ymin=-1,ymax=11] \end{ksys}
$a =$ \_\_
\end{teile}
\end{aufgabe}

%% TODO Titel: Ich-kann-Satz fehlt in der Bank (Platzhalter: Kettenname) – Nr. 17
%% TODO Anweisung: ein Satz über der ersten Teilaufgabe fehlt in der Bank – Nr. 17
\begin{aufgabe}{Parameter aus Bedingungen – Fehler finden, Begründen, Anwendung}
\begin{teile}
\teil Gesucht ist $a > 0$ so, dass der Graph von $f_a(x) = x^2 - a x$ mit der x-Achse die Fläche $\frac{4}{3}$ einschließt. Nils rechnet: $-\frac{a^3}{6} = \frac{4}{3}$, also $a^3 = -8$ und $a = -2$. Finde den Fehler und rechne richtig.
\teil Begründe, warum die Forderung „Der Graph von $f_a$ schließt mit der x-Achse eine Fläche mit dem Inhalt $10$ ein“ eine Gleichung für $a$ liefert.
\teil Ein Beet hat die Form des Flächenstücks zwischen dem Graphen von $b_a(x) = a x \cdot (6 - x)$ mit $a > 0$ und der x-Achse ($x$ und $b_a(x)$ in m). Es soll 18 m² groß sein. Bestimme $a$ und die größte Ausdehnung des Beets in y-Richtung. $a =$ \_\_, Ausdehnung: \_\_ m
\end{teile}
\end{aufgabe}
